%! Properties and Applications of Determinants — Unit 5, Lesson 5.2 — Warm-Up (Answer Key)
\documentclass[10pt]{article}
\usepackage{linalg-article}
\usepackage{linalg-key}   % pulls in -boxes; do NOT also load -boxes
\newcommand{\TallMath}[1]{$\displaystyle #1\rule[-1.4em]{0pt}{3.2em}$}

\begin{document}

\pageheader{Unit 5, Lesson 5.2}{Warm-Up --- Answer Key}
\namedateperiod

\vspace{0.12in}

\begin{notesbox}{Getting Ready: Does Elimination Change the Determinant?}
\small We will run all three steps on the \emph{same} matrix
$M=\begin{bsmallmatrix} 1 & 2 \\ 2 & 5 \end{bsmallmatrix}$. Watch what happens to its determinant.

\vspace{4pt}
\textbf{1. A $2\times2$ determinant (Lesson 5.0).} Recall
$\det\begin{bsmallmatrix} a & b\\ c & d \end{bsmallmatrix}=ad-bc$. Compute
\[
  \det M = \det\begin{bmatrix} 1 & 2 \\ 2 & 5 \end{bmatrix} = {\color{keyred}\mathbf{(1)(5)-(2)(2)=1}}.
\]

\vspace{4pt}
\textbf{2. One elimination step (Unit 2).} Subtract $2\times(\text{row }1)$ from row $2$. Write the new
matrix $U$:
\[
  M=\begin{bmatrix} 1 & 2 \\ 2 & 5 \end{bmatrix}
  \;\xrightarrow{\;\text{row }2\,-\,2\,(\text{row }1)\;}\;
  U=\begin{bmatrix} 1 & 2 \\ {\color{keyred}\mathbf{0}} & {\color{keyred}\mathbf{1}} \end{bmatrix}.
\]

\vspace{4pt}
\textbf{3. A triangular determinant (Lesson 5.1).} Now $U$ is triangular, so its determinant is the
product of the diagonal:
\[
  \det U = {\color{keyred}\mathbf{(1)(1)=1}}.
\]

\vspace{4pt}
\textbf{One-sentence prediction.} Compare $\det M$ and $\det U$. Did the elimination step change the
determinant? \ansline{No --- $\det M=\det U=1$. Adding a multiple of one row to another leaves the determinant unchanged.}
\end{notesbox}

\begin{teachernote}
This is the whole lesson in miniature. Items~1--3 walk one matrix through \emph{compute $\det$},
\emph{one elimination step}, \emph{triangular $\det$} --- and the punchline is that the two
determinants \emph{agree}. Name it: ``adding a multiple of one row to another is a \emph{shear}; it
slides the box but keeps its base and height, so the volume (and the determinant) is unchanged.'' That
single fact is what makes ``$\det=$ product of the pivots'' work in the notes. Watch for $ad+bc$ in
item~1 and for students changing the top row in item~2 (only row~2 changes).
\end{teachernote}

\end{document}
